
% ======================================================================
% ABSTRACT
% ======================================================================
\setstretch{1.5}
\begin{center}
{\fontsize{16}{19}\selectfont\textbf{Abstract}}
\end{center}

\vspace{3mm}

Large language models (LLMs) are developing rapidly, they are opening up new paths for Enterprise Digital Transformation. However, when applying them in some specific  areas, such as banks, huge challenges are needed to be solved. For example, programming code can be wrong, errors happen in some cases, tools may not work well. Also, we depend on outside application programming interfaces (APIs) which brings extra risks. To solve these issues, this article proposes a system and implementation method specifically designed for digital transformation of banks.

The system has three agents. A Code Agent turns natural language questions into executable queries by using a Python-based LLM service and a multi-level SQL whitelist filter. A Retrieval-Augmented Generation (RAG) Agent fetches policies and knowledge by using Milvus vector indexing and vLLM text generation. A Tool Agent automatically handles everyday business tasks, such as spotting schedule conflicts, booking meeting rooms, and calculating routes. These three agents are all managed by a DeepSeek-based module, which understands what the user wants from a central point.

The platform has an access control system based on roles. This system is called RBAC. It gives different levels of security permissions. It also keeps data separate by department. The platform runs on a five-layer microservices architecture. The architecture is built with Spring Boot, Vue 3, and Docker containers.

\newpage
